\documentclass[../main.tex]{subfiles}

\graphicspath{{../image/}{../figure/}}

\begin{document}

\chapter{一维空间的标量守恒律}\label{chapter:ch11-一维标量守恒律}

尽管一维空间的标量守恒律表面上看似简单，却蕴含着极为丰富的理论。它之所以值得关注，不仅在于其内在的理论价值，还在于它能为我们理解方程组的行为提供宝贵的洞见。本章的讨论将运用\cref{chapter:ch10-广义特征线}中发展起来的广义特征线理论。


本章的大部分篇幅将致力于讨论真正非线性的情形，其特点在于：极值向后广义特征线本质上是经典特征线，即解沿之为常数的直线。这一性质施加了极强的约束，使我们得以推导出关于解的正则性与大时间行为的相当精确的信息。

解在近似连续点处是（古典意义下）连续的，并且在连续点集的内部局部 Lipschitz 连续。近似跳跃间断点位于经典激波上。其余的不规则点至多可数，由激波的碰撞和/或压缩波的聚焦所形成。一般地，当初值光滑时，解是分段光滑的。

真正非线性引发了一系列耗散机制，影响着解的大时间行为。熵耗散诱导初值属于 $L^{p}(-\infty,\infty)$ 的解以 $O(t^{-\frac{p}{p+1}})$ 的速率衰减。当初值具有紧支集时，特征线的发散会生成 $N$ 波形剖面。在周期初值下，特征线的被约束诱导出每周期总变差以 $O(t^{-1})$ 的速率衰减，并形成锯齿状剖面。

真正非线性标量守恒律 Cauchy 问题的可容许弱解的另一个重要特征是：它们通过 Lax 函数与其初值显式联系。本章将借助广义特征线方法建立这一性质；它也可以作为发展 Cauchy 问题解的一般理论的另一个出发点。

关于解的比较定理提供了进一步的洞见。下文将证明，任一可容许解的回折数（lap number）随时间单调不增。

此外，任意两个解之间的 $L^{1}$ 距离一般地单调不增，但通常守恒；而适当加权的 $L^{1}$ 距离则严格递减。

广义特征线方法的优点之一在于，它可以方便地推广到非齐次、真正非线性的平衡律。本章将概述这一理论，并给出两个例子，用以展示非齐次项与源项如何影响解的大时间行为。

与之形成对照的是，当真正非线性失效时，我们将面临相当大的复杂性。这将以最简单的、通量函数带有单个拐点的情形为例加以说明。拐点的存在对解的一般正则性与长时间行为具有严重的影响。

广义特征线理论所提供的工具，将被用于判定：在缺失有界变差的情形下，解在连续点处是否可能产生熵。

\subfile{Chapter-11/section-01.tex}
\subfile{Chapter-11/section-02.tex}
\subfile{Chapter-11/section-03.tex}
\subfile{Chapter-11/section-04.tex}
\subfile{Chapter-11/section-05.tex}
\subfile{Chapter-11/section-06.tex}
\subfile{Chapter-11/section-07.tex}
\subfile{Chapter-11/section-08.tex}
\subfile{Chapter-11/section-09.tex}
\subfile{Chapter-11/section-10.tex}
\subfile{Chapter-11/section-11.tex}
\subfile{Chapter-11/section-12.tex}
\subfile{Chapter-11/section-13.tex}
\subfile{Chapter-11/section-14.tex}

\end{document}
